\documentclass[11pt]{article}
\pdfinfoomitdate=1
\pdftrailerid{}
\pdfsuppressptexinfo=15
\usepackage[margin=1in]{geometry}
\usepackage[T1]{fontenc}
\usepackage{lmodern}
\usepackage{microtype}
\usepackage{amsmath,amssymb,amsthm,mathtools}
\usepackage{booktabs,longtable,array}
\usepackage{enumitem}
\usepackage{xcolor}
\usepackage[hidelinks]{hyperref}
\usepackage{fancyhdr}
\usepackage{lastpage}
\newtheorem{theorem}{Theorem}[section]
\newtheorem{lemma}[theorem]{Lemma}
\newtheorem{proposition}[theorem]{Proposition}
\newtheorem{corollary}[theorem]{Corollary}
\newtheorem{definition}[theorem]{Definition}
\newtheorem{remark}[theorem]{Remark}
\newtheorem{firewall}[theorem]{Firewall}
\newcommand{\1}{\mathbf 1}
\newcommand{\RH}{\mathrm{RH}}
\newcommand{\RePart}{\operatorname{Re}}
\pagestyle{fancy}
\fancyhf{}
\lhead{Large-prime parity boundary}
\rhead{\thepage/\pageref{LastPage}}
\setlength{\headheight}{14pt}
\title{\bfseries The Large-Prime Parity Boundary\\[2mm]
\large A No-Go for \texttt{LAPBR67} and an Exact Signed Type-II Bellman Frontier}
\author{Agentic Polymath Project}
\date{18 August 2026}
\begin{document}
\maketitle
\begin{center}
\fcolorbox{red!70!black}{white}{\parbox{0.91\textwidth}{\small
\textbf{Status.} The adaptive small-prime cube is retained.  The proposed
nonnegative large-prime residual is disproved.  The paper replaces it with one
exact largest-prime Bellman identity and one open signed Type-II estimate.
The Riemann Hypothesis is not proved.}}
\end{center}
\begin{abstract}
A recent factor--67 parity proposal proves positivity of the source histories
of length below an adaptive even depth and supported on primes
$p\le (\log X)^{1/4}$.  It asks that the complementary histories be
nonnegative.  We show that this residual theorem, called \texttt{LAPBR67}, is
false.  A product-safe elementary symmetric-sum argument applies at the
published depth $L=O(\log\log\log X)$: the final layer may be restricted to
primes at most $X^{1/(2L)}$, while all preceding layers admit a full reciprocal-
prime upper bound.  Mertens' theorem then makes the last odd layer dominate, so
the complete depth current is
negative, and the residual is more negative than the small-prime cube is
positive.

The repair is to resum every large-prime depth.  We prove positivity of the
complete small-prime cube, assign every remaining history to its largest prime,
and obtain an exact Bellman identity.  The terminal range $p>X/Z$ is
$o(U_Z(X))$.  The only remaining quantity is a one-sided signed bilinear
correlation between the largest prime and a M\"obius tail.  We formulate the
minimal estimate \texttt{BLPTE67}; it implies the annular scalar inequality and
therefore the Riemann Hypothesis through the exact reciprocal-zeta Mellin
consumer.  \texttt{BLPTE67} remains open.
\end{abstract}
\tableofcontents

\section{Genealogy and notation}
Let $b(Y)$ be the complete $P_{61}$ annular $5{:}3$ scalar from the repaired
finite theorem.  It is nonnegative and, once the fixed $P_{61}$ activation
range has been passed,
\begin{equation}\label{eq:base-asymp}
 b(Y)=a_*\sqrt Y+c_*+O(Y^{-3/2}),
 \qquad
 a_*=12\prod_{p\le61}\left(1-\frac1p\right)>0.
\end{equation}
For squarefree rough histories define
\begin{equation}\label{eq:Br}
 B_r(X)=\sum_{\substack{P^-(m)\ge67\\\omega(m)=r}}
 \frac1{\sqrt m}b(X/m),
\end{equation}
with zero extension below support.  The exact depth-$L$ current is
\begin{equation}\label{eq:CL}
 C_L(X)=\sum_{r=0}^{L-1}(-1)^rB_r(X).
\end{equation}


\section{Product-safe domination of the last layer}
Put
\[
 \Lambda_X=\sum_{67\le p\le X}\frac1p=\log\log X+O(1).
\]
The no-go needs no uniform Sathe--Selberg theorem.  It follows from elementary
symmetric sums after restricting only the final layer to a product-safe prime
range.

\begin{lemma}[Product-safe layer comparison]\label{lem:layers}
Let $r=O(\log\log\log X)$, put $Y=X^{1/(2r)}$, and set
\[
 A_Y=\sum_{67\le p\le Y}\frac1p.
\]
Then
\[
 B_r(X)\ge c_0\sqrt X\frac{A_Y^r}{r!}(1-o(1)),
\]
whereas
\[
 \sum_{j<r}B_j(X)
 \le 2C_0\sqrt X\frac{\Lambda_X^{r-1}}{(r-1)!}.
\]
If additionally $r\log r/\Lambda_X\to0$, then
\[
 \frac{B_r(X)}{\sum_{j<r}B_j(X)}\longrightarrow\infty.
\]
\end{lemma}
\begin{proof}
Equation~\eqref{eq:base-asymp} and the finite initial range give constants
$c_0,C_0,Y_0>0$ with
\[
 b(u)\ge c_0\sqrt u\ (u\ge Y_0),\qquad
 0\le b(u)\le C_0\sqrt u\ (u\ge1).
\]
Every product of $r$ primes at most $Y$ is at most $\sqrt X$, so the first
inequality may be used on the entire selected final layer.  Let $e_r(Y)$ be the
elementary symmetric sum of the weights $1/p$, $67\le p\le Y$, and put
$S_2=\sum_{p\ge67}p^{-2}$.  The ordered distinct-prime sum satisfies
\[
 r!e_r(Y)\ge A_Y^r-\binom r2S_2A_Y^{r-2};
\]
this is the union bound for ordered tuples with a repeated coordinate.  Since
$r/A_Y\to0$, this proves the lower bound.

For $j<r$, the global upper bound gives
$B_j(X)\le C_0\sqrt X\Lambda_X^j/j!$.  Since $r/\Lambda_X\to0$, these terms
increase throughout the relevant range and their sum is at most twice the last
one.

Finally Mertens' theorem yields
\[
 A_Y=\Lambda_X-\log(2r)+O(1).
\]
Consequently
\[
 \frac{A_Y}{r}\left(\frac{A_Y}{\Lambda_X}\right)^{r-1}\to\infty
\]
because $r\log r/\Lambda_X\to0$.
\end{proof}

\section{The published residual is negative}
Set
\[
 Z=(\log X)^{1/4},\qquad z=\sum_{67\le p\le Z}\frac1p,
\]
and let $L$ be the least even integer with $L\ge8(z+1)$.  Then
\[
 L=O(\log\log\log X),\qquad L/\Lambda_X\to0.
\]

\begin{theorem}[No-go for \texttt{LAPBR67}]\label{thm:no-go}
For all sufficiently large $X$,
\[
 C_L(X)<0.
\]
If $S_X$ is the positive small-prime cube of the predecessor and
$\mathcal L_X=C_L(X)-S_X$ is its large-prime residual, then
\[
 \boxed{\mathcal L_X<-S_X<0.}
\]
Moreover $-\mathcal L_X/S_X\to\infty$.
\end{theorem}
\begin{proof}
Apply Lemma~\ref{lem:layers} with $r=L-1$.  Here
$r=O(\log\log\log X)$ and $r\log r/\Lambda_X\to0$.  The final term of
\eqref{eq:CL} is negative because $L$ is even, and its magnitude exceeds the
sum of every preceding positive layer.  Thus $C_L(X)<0$.  The residual
inequality follows from $S_X>0$.  Moreover
$S_X\ll\sqrt X e^z$, while the product-safe lower bound for $B_{L-1}$ has
logarithm $(L-1)\log(\Lambda_X/(L-1))+O(L)$ after division by $\sqrt X$;
this dominates $z$, so the ratio tends to infinity.
\end{proof}

\begin{firewall}
The theorem does not refute the small-prime cube.  It refutes the stopping rule
that asks all histories meeting a larger prime to form a nonnegative truncated
current.
\end{firewall}

\section{The complete small-prime cube}
Let $Q_Z=\prod_{67\le p\le Z}p$ and define
\[
 \mathcal B_Z(X)=\sum_{d\mid Q_Z}\frac{\mu(d)}{\sqrt d}b(X/d),
 \qquad U_Z(X)=\mathcal B_Z(X)/\sqrt X.
\]

\begin{theorem}[Complete cube positivity]\label{thm:small-complete}
For all sufficiently large $X$,
\begin{align}
 \mathcal B_Z(X)
 ={}&a_*\sqrt X\prod_{67\le p\le Z}(1-p^{-1})\notag\\
 &+c_*\prod_{67\le p\le Z}(1-p^{-1/2})+o(1)>0.
\end{align}
In particular $U_Z(X)\asymp1/\log Z$.
\end{theorem}
\begin{proof}
The largest divisor of $Q_Z$ is $\exp((1+o(1))Z)=X^{o(1)}$, so the asymptotic
\eqref{eq:base-asymp} is uniform over every divisor.  The main and constant
terms factor exactly.  The remainder is bounded by
\[
 X^{-3/2}\sum_{d\mid Q_Z}d
 =X^{-3/2}\prod_{67\le p\le Z}(1+p)=o(1).
\]
Mertens' theorem finishes the claim.
\end{proof}

\section{Largest-prime Bellman ownership}
Adjoin active primes $Z<p_1<\cdots<p_k\le X/2$ in increasing order.  The exact
normalized finite-prime recurrence is
\[
 U_i(Y)=U_{i-1}(Y)-p_i^{-1}U_{i-1}(Y/p_i),\qquad U_0=U_Z.
\]

\begin{theorem}[Largest-prime Bellman identity]\label{thm:bellman}
The complete normalized scalar is
\begin{equation}\label{eq:bellman}
 U_{\rm full}(X)=U_Z(X)-
 \sum_{Z<p\le X/2}\frac1pU_{<p}(X/p).
\end{equation}
Moreover
\begin{equation}\label{eq:child-expansion}
 U_{<p}(Y)=
 \sum_{\substack{v\ \mathrm{squarefree}\\Z<P^-(v),\ P^+(v)<p}}
 \frac{\mu(v)}vU_Z(Y/v).
\end{equation}
\end{theorem}
\begin{proof}
Telescope the recurrence.  Expanding the previously adjoined primes in each
child yields~\eqref{eq:child-expansion}.  Every nonempty subset occurs exactly
once, in the summand indexed by its largest prime.
\end{proof}

The predecessor's truncated residual also has the exact Duhamel identity
\[
 P_{<L}(R)b-P_{<L}(S)b
 =-\sum_{a+b\le L-2}(-1)^{a+b}S^a(R-S)R^b b,
\]
which displays its short remaining histories.  Theorem~\ref{thm:no-go} shows
why they must be resummed rather than signed separately.

\section{Type I and Type II}
Split the Bellman budget in~\eqref{eq:bellman} at $p=X/Z$:
\[
 \mathfrak I_Z(X)=\sum_{X/Z<p\le X/2}\frac1pU_Z(X/p),
\]
\[
 \mathfrak T_Z(X)=\sum_{Z<p\le X/Z}\frac1pU_{<p}(X/p).
\]
The Type-I equality is exact because $X/p<Z$, so no prime above $Z$ is active.

\begin{proposition}[Terminal Type-I closure]\label{prop:typei}
\[
 |\mathfrak I_Z(X)|\ll\frac{(\log Z)^2}{\log X}=o(U_Z(X)).
\]
\end{proposition}
\begin{proof}
For $2\le y\le Z$, unsigned divisor mass gives
$|U_Z(y)|\ll\log(2y)$.  The reciprocal-prime mass of
$(X/Z,X/2]$ is $O(\log Z/\log X)$ by Mertens' theorem.
\end{proof}

The remaining term is the exact signed Type-II correlation
\begin{equation}\label{eq:typeii}
 \mathfrak T_Z(X)=
 \sum_{Z<p\le X/Z}\frac1p
 \sum_{\substack{v\ \mathrm{squarefree}\\Z<P^-(v),\ P^+(v)<p}}
 \frac{\mu(v)}vU_Z(X/(pv)).
\end{equation}
No absolute values have been taken in~\eqref{eq:typeii}.

\begin{definition}[\texttt{BLPTE67}]
The balanced large-prime Type-II estimate is
\[
 \boxed{\mathfrak T_Z(X)\le U_Z(X)-\mathfrak I_Z(X)}
\]
at every sufficiently large admissible state.
\end{definition}

\section{Conditional completion}
\begin{theorem}[Exact remaining theorem]\label{thm:conditional}
\[
 \texttt{BLPTE67}\Longrightarrow\RH.
\]
\end{theorem}
\begin{proof}
The Bellman identity and \texttt{BLPTE67} give nonnegativity at each state.
Child endpoints are smaller, so finite induction starts from the directly
checked base range.  At the root one obtains eventual nonnegativity of
\[
 \mathcal A_X=5[c_X(2)-c_{X/4}(2)]+3[c_X(3)-c_{X/4}(3)].
\]
Its Mellin transform is
\[
 (1-4^{-s})\left[\frac6{s^2}-
 \frac{3(1-2^{-z})(2-2^{-z})}{s^2\zeta(z)}\right],
 \qquad z=s+\tfrac12.
\]
The elementary factors do not cancel a zeta zero in $\Re z>0$.  Landau's
theorem excludes zeros to the right of the critical line, and the functional
equation supplies the reflection.
\end{proof}

\section{Proof boundary}
\begin{longtable}{p{0.42\textwidth}p{0.42\textwidth}}
\toprule
Statement & Status\\
\midrule
Adaptive small-prime cube & retained\\
\texttt{LAPBR67} residual nonnegativity & false\\
Product-safe last-layer domination & proved from symmetric sums and Mertens\\
Complete small-prime cube & proved positive\\
Largest-prime Bellman identity & exact\\
Terminal Type I & closed\\
Signed balanced Type II, \texttt{BLPTE67} & open / RH-bearing\\
Riemann Hypothesis & unproved\\
\bottomrule
\end{longtable}

\section{Conclusion}
The large prime does simplify the source, but not by making a shallow residual
positive.  It supplies a unique owner, a strict endpoint reduction, a terminal
Type-I range, and a canonical prime-versus-M\"obius Type-II form.  The original
residual sign is false.  The exact remaining theorem is the one-sided estimate
\texttt{BLPTE67}; proving it, without discarding the M\"obius signs, would finish
this lane.
\end{document}
